% Glossary additions identified from Chapter 4 (Classical vs Quantum Computing)

\newglossaryentry{bit}{
    name={bit},
    description={The fundamental unit of information in classical computing, representing either a 0 or a 1. Contrasted with \gls{qubit}. See Section \ref{sec:differences_rev}.}
}

\newglossaryentry{probability-amplitude}{
    name={probability amplitude},
    description={A complex number whose squared magnitude represents the probability of measuring a quantum system (like a \gls{qubit}) in a specific basis state. Used in the description of \gls{superposition}. See Section \ref{sec:differences_rev}.}
}

\newglossaryentry{computational-complexity-theory}{
    name={computational complexity theory},
    description={The branch of theoretical computer science that studies the resources (e.g., time, memory) required to solve computational problems, classifying them into complexity classes like \gls{p-complexity}, \gls{np-complexity}, \gls{bpp-complexity}, and \gls{bqp}. See Section \ref{sec:complexity_classes_rev}.}
}

\newglossaryentry{turing-machine}{
    name={Turing machine},
    description={A mathematical model of computation that defines an abstract machine manipulating symbols on a strip of tape according to rules. A fundamental model for classical computation. Mentioned in Section \ref{subsec:classical_rev}.}
}

\newglossaryentry{von-neumann-architecture}{
    name={von Neumann architecture},
    description={A computer architecture based on the concept of stored-program computers where instruction data and program data are stored in the same memory. The basis for most modern classical computers. Mentioned in Section \ref{subsec:classical_rev}.}
}

\newglossaryentry{classical-core}{
    name={classical core},
    description={A standard processing unit within a classical computer's central processing unit (CPU) that executes instructions sequentially or in parallel with other cores. Mentioned in Section \ref{subsec:classical_rev} in the context of parallel processing.}
} 